% !PRINTED 10
\weberpage{10}

%<<CONT
dass $\alpha$ der kleinste positive Exponent sei, für den $a^\alpha = 1$
ist. Es ist immer und nur dann
\[
a^\mu = a^{\mu'},
\]
wenn $\mu \equiv \mu'\ (\mathrm{mod}\ \alpha)$, und hierin können die
Exponenten auch negativ sein, wenn $a^{-r}$ die $r^{\mathrm{te}}$ Potenz
des Elementes $a^{-1}$ oder das zu $a^r$ entgegengesetzte Element
bedeutet.

Die Reihe
\[
A = 1,\ a,\ a^2 \ldots a^{\alpha-1}
\]
enthält dann $\alpha$ von einander verschiedene Elemente von $P$, wobei
alle Potenzen von $a$ vertreten sind. Die Elemente $A$ bilden aber
offenbar eine Gruppe vom Grade $\alpha$, weil sich die Exponenten bei der
Zusammensetzung einfach addiren. Diese Gruppe ist ein Theiler von $P$ und
also ist $\alpha$ \gesperrt{ein Theiler von $n$}. Es ist also für jedes Element
$a^n = 1$.

\gesperrt{Die Gruppe $A$ heisst die Periode des Elementes $a$ und $\alpha$ wird auch
der Grad des Elementes $a$ genannt.}

\begin{center}
\S.\,3.

\gesperrt{Normaltheiler einer Gruppe.}
\end{center}

Ist $P$ wie oben eine Gruppe und $Q$ ein Divisor von $P$, ferner $b$ ein
nicht in $Q$ enthaltenes Element von $P$, so ist die Nebengruppe $Qb$
keine Gruppe. Dagegen bildet das System $b^{-1}Qb$, oder ausführlich
geschrieben
\[
1,\ b^{-1}a_1 b,\ b^{-1}a_2 b \ldots b^{-1}a_{\nu-1} b
\]
sicher eine Gruppe, weil
\[
b^{-1}a_1 b\ .\ b^{-1}a_2 b = b^{-1}a_1 a_2 b
\]
ist, und diese Gruppe ist mit $Q$ isomorph. Die Gruppe $b^{-1}Qb$ heisst
\gesperrt{die durch $b$ transformirte Gruppe}. Gehört $b$ selbst zu $Q$, so ist
$b^{-1}Qb$ mit $Q$ identisch. Nimmt man für $b$ die verschiedenen Elemente
von $P$, so erhält man eine ganze Schaar solcher Gruppen, die wir die
\gesperrt{zu $Q$ conjugirten Theiler von $P$} oder auch kurz \gesperrt{conjugirte Gruppen} nennen.

Es kann vorkommen, dass alle conjugirten Theiler mit einander identisch
sind. Wir haben schon bei den Permutationsgruppen gesehen, dass dieser
Fall von besonderer Wichtigkeit ist, und führen also nun allgemein
folgende Definition ein:
%CONT>>
